\chapter{مبانی و پیشینه}
\label{ch:background}

\section{رادار روزنه مصنوعی و \lr{Sentinel-1}}
رادار روزنه مصنوعی (\lr{SAR}) پالس‌های ریزموج به سطح زمین می‌فرستد و بازتاب آن‌ها را ثبت
می‌کند؛ با پردازش سیگنال در طول مسیر ماهواره، تفکیکی به دست می‌آید که با آنتنی واقعی
ناممکن بود. سطح آرام دریا بیشتر انرژی را دور از ماهواره بازمی‌تاباند و تیره دیده می‌شود، اما
بدنه و سازه فولادی شناور مانند بازتابنده‌ای گوشه‌ای عمل می‌کند و نقطه‌ای روشن پدید می‌آورد.

مأموریت \lr{Sentinel-1} \cite{torres2012} تصاویر باند \lr{C} را در حالت پهنای گسترده
(\lr{IW}) با نوار حدود ۲۵۰ کیلومتر برمی‌دارد. در این پژوهش محصول \lr{GRD} با پیکسل ۱۰ متر و
دو قطبش \lr{VV} و \lr{VH} به کار رفت؛ شناور در قطبش متقاطع \lr{VH} معمولاً تضاد بیشتری با دریا
دارد. تصاویر از زیرساخت \lr{Copernicus Data Space Ecosystem} \cite{cdse} و به قالب
\lr{Cloud-Optimized GeoTIFF} خوانده شد، که خواندن تنها پنجره مورد نیاز را بدون بارگیری کل
تصویر ممکن می‌کند. در دوره این پژوهش ماهواره‌های \lr{Sentinel-1C} و \lr{Sentinel-1D} جای
\lr{Sentinel-1A} را گرفتند؛ ازاین‌رو تعمیم مدل به این ماهواره‌ها جداگانه سنجیده شد.

\section{سامانه شناسایی خودکار و \lr{Global Fishing Watch}}
\lr{AIS} برای ایمنی ناوبری طراحی شد و شناور شناسه (\lr{MMSI})، موقعیت، سرعت و جهتش را
پخش می‌کند. پوشش، صداقت و پیوستگی آن تضمین‌شده نیست: فرستنده ممکن است خاموش شود و
موقعیت‌ها ممکن است جعل یا با اخلال \lr{GNSS} جابه‌جا شوند. \lr{Global Fishing Watch} داده‌های
\lr{AIS} را در مقیاس جهانی پردازش می‌کند \cite{kroodsma2018} و آشکارسازی‌های شناور در تصاویر
\lr{Sentinel-1} را با \lr{AIS} تطبیق می‌دهد \cite{paolo2024}. این پژوهش از دو داده عمومی آن
استفاده کرد \cite{gfwapi}: آشکارسازی‌های راداری (با و بدون تطبیق \lr{AIS}) در خانه‌های
\LRE{0٫01} درجه‌ای (حدود ۱ کیلومتر) و ساعتی، و حضور ساعتی شناورهای \lr{AIS} در همان خانه‌ها.
\lr{GFW} سازه‌های ثابت را جدا می‌کند و در نوار نزدیک ساحل آشکارسازی گزارش نمی‌کند؛ این دو
نکته در فصل نتایج اهمیت می‌یابند.

\section{آشکارسازی کلاسیک: \lr{CFAR}}
روش کلاسیک آشکارسازی شناور در \lr{SAR}، آشکارساز با نرخ هشدار نادرست ثابت (\lr{CFAR}) است
\cite{crisp2004}. برای هر پیکسل، میانگین بازتاب دریا در حلقه‌ای از پیکسل‌های آموزشی پیرامون آن
برآورد می‌شود (با نواری محافظ تا انرژی خود هدف برآورد را آلوده نکند) و پیکسلی که از مضربی از
این میانگین روشن‌تر باشد آشکار شمرده می‌شود. مضرب از آماره دریا و نرخ هشدار نادرست
دلخواه به دست می‌آید. \lr{CFAR} به داده آموزشی نیاز ندارد و پارامترهایش تفسیرپذیرند، پس
خط مبنای مناسبی است؛ اما بازتاب سازه‌ها، ساحل و نوارهای اثر رادار را نیز آشکار می‌کند و
شناورهای بزرگ را گاه به چند تکه می‌شکند.

\section{یادگیری عمیق برای آشکارسازی شناور}
مسابقه \lr{xView3} \cite{paolo2022} داده‌ای بزرگ از تصاویر \lr{Sentinel-1} با برچسب شناور و
وضعیت \lr{AIS} فراهم کرد. راه‌حل برنده \cite{xview3first} شناور را به‌جای جعبه مرزی، به‌صورت
نقطه آشکار می‌کرد: شبکه‌ای با رمزگذار از پیش آموزش‌دیده و رمزگشای \lr{U-Net}
\cite{ronneberger2015} برای هر پیکسل خروجی، احتمال مرکز شناور و جابه‌جایی زیرپیکسلی آن را
پیش‌بینی می‌کند، به روش «اشیا به‌مثابه نقطه» \cite{zhou2019}. این رویکرد برای اهدافی که در
تصویر تنها چند پیکسل‌اند مناسب‌تر از آشکارسازهای جعبه‌ای است. رمزگذار این پژوهش
\lr{EfficientNet-B0} \cite{tan2019} است و تابع زیان نقشه مرکز، زیان کانونی کاهش‌یافته
\cite{lin2017,zhou2019}.

\section{شناورهای تاریک و تلفیق داده‌ها}
یافتن شناور تاریک یعنی پیدا کردن آشکارسازی راداری‌ای که هیچ شناور \lr{AIS} در آن زمان و
مکان آن را توضیح ندهد. دشواری در جزئیات است: زمان گذر ماهواره و پیام‌های \lr{AIS} یکی نیست،
موقعیت‌ها خطا دارند، شناور میان پیام‌ها جابه‌جا می‌شود و چند شناور ممکن است به یک
آشکارسازی نزدیک باشند. این پژوهش تطبیق را به‌صورت تخصیص یک‌به‌یک
\cite{kuhn1955} میان آشکارسازی‌ها و «دروازه‌های» زمانی-مکانی شناورهای \lr{AIS} انجام داد
(بخش~\ref{sec:fusion}).
